\documentclass[a4paper,9pt]{extarticle}

%=========================================================
% PACKAGES
%=========================================================

\usepackage[
a4paper,
top=1.8cm,
bottom=1.5cm,
footskip=7mm,
left=1cm,
right=1cm,
headheight=14pt,
headsep=5mm
]{geometry}
\usepackage[T1]{fontenc}
\usepackage[utf8]{inputenc}
\usepackage[french]{babel}

\usepackage[sfdefault]{FiraSans}
\renewcommand{\familydefault}{\sfdefault}

\usepackage{multicol}
\setlength{\columnsep}{0.7cm}

\usepackage{amsmath,amssymb}

\usepackage{titlesec}
\usepackage{enumitem}

\usepackage{tcolorbox}
\tcbuselibrary{skins,breakable}

\usepackage{xcolor}

\usepackage{graphicx}
\usepackage{float}

\usepackage{tikz}

\usepackage{fancyhdr}
\usepackage[hidelinks]{hyperref}


%=========================================================
% PAGE SETTINGS
%=========================================================

\setlength{\parindent}{0pt}
\setlength{\parskip}{1.5pt}

\setlength{\abovedisplayskip}{4pt}
\setlength{\belowdisplayskip}{4pt}

\setlist[itemize]{
leftmargin=*,
itemsep=1pt,
topsep=1pt
}


%=========================================================
% COLORS
%=========================================================

\definecolor{Blue}{HTML}{E3F0FF}
\definecolor{Green}{HTML}{E3F7EA}
\definecolor{Pink}{HTML}{FFE5ED}
\definecolor{Orange}{HTML}{FFF0DC}
\definecolor{Purple}{HTML}{EEE7FF}
\definecolor{Yellow}{HTML}{FFFADB}


\definecolor{BlueTitle}{HTML}{285B8F}
\definecolor{GreenTitle}{HTML}{287052}
\definecolor{PinkTitle}{HTML}{A84063}
\definecolor{OrangeTitle}{HTML}{A96500}
\definecolor{PurpleTitle}{HTML}{6245A5}


%=========================================================
% HEADER
%=========================================================

\pagestyle{fancy}

\fancyhf{}

\fancyhead[L]{\small\textbf{Computer Security}}
\fancyhead[C]{\small Security Design Principles}
\fancyhead[R]{\small Week 02 --- 14-20/09/2026}

\fancyfoot[C]{\small\href{https://github.com/AlphaXT38G/notes_ba5}{github.com/AlphaXT38G/notes\_ba5}}
\fancyfoot[R]{\small \thepage}

\renewcommand{\headrulewidth}{0.4pt}



%=========================================================
% TITLES
%=========================================================

\titleformat{\section}
{\large\bfseries\color{BlueTitle}}
{}{0pt}{}

\titlespacing{\section}
{0pt}{6pt}{3pt}


% separator

\newcommand{\separator}{
\vspace{2pt}
\hrule
\vspace{3pt}
}


%=========================================================
% BOXES
%=========================================================

% DEFINITION - blue
\newtcolorbox{definition}{
enhanced,
breakable,
colback=Blue,
colframe=BlueTitle,
boxrule=0.7pt,
arc=0mm,
title=\textbf{Definition},
fonttitle=\bfseries
}

% IMPORTANT - pink
\newtcolorbox{important}{
enhanced,
breakable,
colback=Pink,
colframe=PinkTitle,
boxrule=0.7pt,
arc=0mm,
title=\textbf{Key idea},
fonttitle=\bfseries
}

% FORMULAS - green
\newtcolorbox{formula}{
enhanced,
breakable,
colback=Green,
colframe=GreenTitle,
boxrule=0.7pt,
arc=0mm,
title=\textbf{Key formulas},
fonttitle=\bfseries
}

% WARNING - orange
\newtcolorbox{warning}{
enhanced,
breakable,
colback=Orange,
colframe=OrangeTitle,
boxrule=0.7pt,
arc=0mm,
title=\textbf{Warning},
fonttitle=\bfseries
}

% EXAMPLE - purple
\newtcolorbox{example}{
enhanced,
breakable,
colback=Purple,
colframe=PurpleTitle,
boxrule=0.7pt,
arc=0mm,
title=\textbf{Example},
fonttitle=\bfseries
}

% SUMMARY - orange
\newtcolorbox{summary}{
enhanced,
breakable,
colback=Yellow,
colframe=OrangeTitle,
boxrule=0.7pt,
arc=0mm,
title=\textbf{Summary},
fonttitle=\bfseries
}



%=========================================================
% FIGURE BOX
%=========================================================

\newtcolorbox{figurebox}{
enhanced,
colback=white,
colframe=gray,
boxrule=0.5pt,
arc=0mm,
center
}



%=========================================================
% DOCUMENT
%=========================================================

\begin{document}

%\footnotesize   --> To reduce the size of the text

\thispagestyle{fancy}



\begin{multicols}{2}


%=========================================================
\section{Heartbleed bug}


\begin{example}

\textbf{Heartbleed bug:} Issue: The attacker could lie about the size of their data.

\begin{itemize}

\item Attacker sends 1 byte of data.

\item Attacker claims they sent 64,000 bytes.

\item The server, without checking the actual length, would reply with the attacker's 1
byte plus the next 63,999 bytes of whatever was in its private memory.

\end{itemize}

\end{example}


\begin{definition}

\textbf{Transport Layer Security (TLS)} is a cryptographic protocol that provides privacy, data integrity, and authentication for internet communications.

\end{definition}



%=========================================================
\section{1 -- Economy of mechanism}


\begin{important}

``Keep the [security mechanism / implementation] design as simple and small as possible''

\end{important}


\textbf{Why?}

It needs to be easy to audit and verify.

(operational testing is not appropriate to evaluate security)

[Penetration testing is valuable]


\separator


\begin{definition}

\textbf{The ``Trusted Computing Base'' (TCB)}

Every component of the system on which the security policy relies.

Hardware / firmware / software

\end{definition}


\begin{important}

The TCB is trusted to operate correctly for the security policy to hold

\end{important}


\begin{warning}

The only proper use of the verb ``to trust'' in Security Engineering:

``X trusts Y will do Z''

\end{warning}



%=========================================================
\section{2 -- Fail-safe defaults}


\begin{important}

``Base access decisions on permission rather than exclusion'' [SS75]

\end{important}


\begin{definition}

\textbf{Ransomware -- What is it?}

\begin{itemize}

\item A Ransomware denies you access to your own data, most commonly by encrypting your files

\item Attacker then demands a ransom payment in exchange for the private key required to decrypt
your files

\end{itemize}

\end{definition}



\columnbreak

%=========================================================
\section{3 -- Complete mediation}


\begin{important}

``Every access to every object must be checked for authority'' [SS75]

\end{important}



%=========================================================
\section{4 -- Open design}


\begin{important}

``The design should not be secret'' [SS75]

\end{important}



%=========================================================
\section{5 -- Separation of privilege}


\begin{important}

``No single accident, deception, or breach of trust is sufficient to
compromise the protected information'' [SS75]

\end{important}



%=========================================================
\section{6 -- Least privilege}


\begin{important}

``Every program and every user of the system should operate using the
least set of privileges necessary to complete the job'' [SS75]

\end{important}



%=========================================================
\section{7 -- Least common mechanism}


\begin{important}

``Minimize the amount of mechanism common to more than one
user and depended on by all users'' [SS75]

\end{important}



%=========================================================
\section{8 -- Psychological acceptability}


\begin{important}

``It is essential that the human interface be designed for ease of use, so that users
routinely and automatically apply the protection mechanisms correctly'' [SS75]

\end{important}


\separator


\textbf{Extra principles from physical security}



%=========================================================
\section{9 -- Work factor}


\begin{important}

``Compare the cost of circumventing the mechanism with
the resources of a potential attacker'' [SS75]

\end{important}



%=========================================================
\section{10 -- Compromise recording}


\begin{important}

``Reliably record that a compromise of information has occurred [...] in
place of more elaborate mechanisms that completely prevent loss'' [SS75]

\end{important}



\end{multicols}


\end{document}
